\begin{abstract}

Model-heterogeneous federated learning (HFL) allows clients with different architectures to collaborate. Existing corruption-robust pipelines, however, usually treat corruption as label-preserving nuisance variation. We study \textbf{class--corruption entanglement (CLE)}, a controlled failure mode in which client-specific associations between task classes and corruption operators turn corruption into a directional class cue. To distinguish this behavior from ordinary image degradation, we evaluate the same semantic source under multiple operators and introduce \textbf{Directional Shortcut Alignment (DSA)}, a probability-scale contrast aligned with the pre-registered training binding. We characterize the assumptions under which DSA identifies this paired operator response and prove that improvements in an unconstrained pseudo-environment proxy cannot, by themselves, certify reduced CLE dependence. Across three matched training seeds, the Local CLE contrast closely tracks the HFL contrast, while the communication add-on changes sign, supporting predominantly local-first formation. The resulting intervention is communication-agnostic: a taxonomy-assisted coarse \textbf{Public Environment Witness (PEW)} supplies pseudo-environments, and \textbf{Balanced Environment Risk (BER)} compresses class-conditional support imbalance across them. Completed controlled studies show DSA reductions across binding maps and communication bases. Under the seed-0 main evaluation protocol, DSA rises monotonically with controlled binding strength; in a batch-matched AsymHFL comparison, PEW+BER reduces DSA by 51.21\% on all four clients while incurring a nonuniform clean-utility cost. The method is therefore a targeted mitigation for controlled CLE, not a taxonomy-free or universally utility-preserving HFL plug-in.
\ifdraftskeleton
\evidenceneeded{replace the final empirical summary after resolving the remaining main-evaluation Axis-II, second-task, and taxonomy-stress evidence slots.}
\fi

\end{abstract}
